\section{Conclusion}

Two-factor linear networks provide an exact finite-width closure that is broader than the special regimes in which their learning dynamics are usually solved.  For arbitrary finite data, differentiable loss, hidden width, and input dimension, the parameter flow projects to $(f,K)$, independently of the balancing leaf.  The resulting law is driven rather than kernel-autonomous, and exact elimination of $K$ turns it into a full-history Volterra equation for $f$.  The same dynamics are therefore Markovian or memory-bearing according to the observable retained.

The scalar deep-linear hierarchy pinpoints the first rupture.  From three factors onward, generic fibers with fixed output and fixed tangent kernel carry varying $e_{d-2}$ and hence varying kernel velocity.  This failure is not irreducible memory: balancing labels restore closure, and at depth three the required invariant coordinate completes the symmetric hierarchy.  The transition is thus from universal leaf-independent closure to leaf-conditioned closure.  Whether analogous reductions hold on training-reachable sets of nonlinear networks is a separate empirical question.
